\documentclass{article}
\usepackage{amsmath,amssymb,amsthm}
\usepackage{algorithm,algorithmic}
\usepackage{fullpage}
\usepackage{hyperref}
\newtheorem{theorem}{Theorem}
\newtheorem{lemma}{Lemma}
\newcommand{\E}{\mathbb{E}}
\newcommand{\Var}{\operatorname{Var}}
\newcommand{\norm}[1]{\left\|#1\right\|}
\newcommand{\grad}{\nabla}
\begin{document}

\section*{Algorithm Name}
\textbf{Differentially Private Stochastic Gradient Descent with Gaussian Noise (DP‑SGD).}

\section*{Mathematical Formulation}
Let $f:\mathbb{R}^{d}\rightarrow\mathbb{R}$ be the (possibly non‑convex) loss function.  
At iteration $t\in\{0,1,\dots,T-1\}$ we have the iterate $x_t\in\mathbb{R}^{d}$.
The algorithm draws a (mini‑batch) stochastic gradient $g_t$ and adds isotropic
Gaussian noise:
\[
\boxed{
\begin{aligned}
g_t &:= \grad f(x_t;\xi_t) + \zeta_t,\\[0.2em]
\zeta_t &\sim \mathcal{N}\!\bigl(0,\sigma^{2} I_d\bigr),\\[0.2em]
x_{t+1} &:= x_t - \eta_t\, g_t .
\end{aligned}}
\]
Here
\begin{itemize}
    \item $\eta_t>0$ is the learning‑rate (possibly time‑varying);
    \item $\xi_t$ denotes the random data sample (or mini‑batch) used to compute the
          stochastic gradient $\grad f(x_t;\xi_t)$;
    \item $\sigma^{2}$ is the variance of the added Gaussian noise, which is
          calibrated to a prescribed privacy budget $(\varepsilon,\delta)$ via the
          moments accountant or R\'enyi DP.
\end{itemize}

\section*{Pseudocode}
\begin{algorithm}[H]
\caption{DP‑SGD}
\begin{algorithmic}[1]
\REQUIRE Initial point $x_0\in\mathbb{R}^{d}$, learning‑rate schedule $\{\eta_t\}_{t=0}^{T-1}$,
         noise scale $\sigma$, clipping norm $C$ (optional), total iterations $T$.
\FOR{$t=0$ \TO $T-1$}
    \STATE Sample a mini‑batch $\mathcal{B}_t$ from the dataset.
    \STATE Compute per‑example gradients $g_{t,i}= \grad f(x_t;z_i)$ for $z_i\in\mathcal{B}_t$.
    \STATE (Optional) Clip each $g_{t,i}$: $ \tilde g_{t,i}= g_{t,i}\cdot\min\!\bigl\{1,\frac{C}{\norm{g_{t,i}}}\bigr\}$.
    \STATE Form the averaged (clipped) gradient $\bar g_t=\frac{1}{|\mathcal{B}_t|}\sum_{i\in\mathcal{B}_t}\tilde g_{t,i}$.
    \STATE Sample noise $\zeta_t\sim\mathcal{N}(0,\sigma^{2}I_d)$.
    \STATE Set $g_t:=\bar g_t+\zeta_t$.
    \STATE Update $x_{t+1}:=x_t-\eta_t g_t$.
\ENDFOR
\RETURN $x_T$ (or the averaged iterate $\bar x_T=\frac{1}{T}\sum_{t=0}^{T-1}x_t$).
\end{algorithmic}
\end{algorithm}

\section*{Dry Run Example}
Consider the simple scalar quadratic $f(w)=(w-3)^2$.
\begin{itemize}
    \item Gradient: $\grad f(w)=2(w-3)$.
    \item Choose $\eta=0.1$, noise variance $\sigma^{2}=0.01$, and no stochasticity
          (i.e. $\xi_t$ is the whole data, so $\grad f(w;\xi_t)=\grad f(w)$).
    \item Initialise $w_0=0$.
\end{itemize}
We perform three iterations.

\begin{center}
\begin{tabular}{c|c|c|c}
$t$ & $g_t= \grad f(w_t)+\zeta_t$ & $w_{t+1}=w_t-\eta g_t$ & $f(w_{t+1})$\\\hline
0 & $2(0-3)+\zeta_0 = -6+\zeta_0$ & $0-0.1(-6+\zeta_0)=0.6-0.1\zeta_0$ & $(0.6-0.1\zeta_0-3)^2$\\
1 & $2(w_1-3)+\zeta_1$ & $w_1-0.1 g_1$ & $f(w_2)$\\
2 & $2(w_2-3)+\zeta_2$ & $w_2-0.1 g_2$ & $f(w_3)$\\
\end{tabular}
\end{center}
If we draw a concrete noise sample, say $\zeta_0=0.2$, $\zeta_1=-0.1$, $\zeta_2=0.05$,
the numerical values are:

\begin{itemize}
    \item $t=0$: $g_0=-5.8$, $w_1=0.58$, $f(w_1)=5.80$.
    \item $t=1$: $\grad f(w_1)=2(0.58-3)=-4.84$, $g_1=-4.94$, $w_2=1.074$, $f(w_2)=3.71$.
    \item $t=2$: $\grad f(w_2)=2(1.074-3)=-3.852$, $g_2=-3.802$, $w_3=1.454$, $f(w_3)=2.41$.
\end{itemize}
The objective value decreases despite the injected noise, illustrating a typical DP‑SGD trajectory.

\newpage
\section*{Assumptions}
\begin{enumerate}
    \item[\textbf{A1}] \textbf{Smoothness.} $f$ is $L$‑smooth:
    \[
        \forall x,y\in\mathbb{R}^{d},\qquad
        f(y)\le f(x)+\langle\grad f(x),y-x\rangle+\frac{L}{2}\norm{y-x}^{2}.
    \tag{A1}
    \]
    \item[\textbf{A2}] \textbf{Bounded variance of stochastic gradients.} There exists $G^{2}>0$ such that
    \[
        \E_{\xi_t}\!\bigl[\norm{\grad f(x_t;\xi_t)-\grad f(x_t)}^{2}\mid x_t\bigr]\le G^{2},
        \qquad\forall t.
    \tag{A2}
    \]
    \item[\textbf{A3}] \textbf{Gradient norm bound.} For all $x$, $\norm{\grad f(x)}\le B$ (a constant $B$ may be infinite; the bound is used only where explicitly needed).
    \item[\textbf{A4}] \textbf{Noise independence.} The added Gaussian noises $\{\zeta_t\}$ are i.i.d.,
    $\zeta_t\sim\mathcal{N}(0,\sigma^{2}I_d)$ and independent of $\xi_t$.
    \item[\textbf{A5}] \textbf{Learning‑rate schedule.} $\eta_t\equiv\eta$ is constant and satisfies $0<\eta\le \frac{1}{L}$.
\end{enumerate}

\section*{Main Convergence Theorem}
\begin{theorem}[Expected Gradient Norm Decrease for DP‑SGD]\label{thm:main}
Assume \textbf{A1}–\textbf{A5} hold. Let $x_0$ be arbitrary and define the averaged iterate
\[
\bar x_T:=\frac{1}{T}\sum_{t=0}^{T-1}x_t .
\]
Then for any $T\ge1$,
\[
\frac{1}{T}\sum_{t=0}^{T-1}\E\!\bigl[\norm{\grad f(x_t)}^{2}\bigr]
\;\le\;
\underbrace{\frac{2\bigl(f(x_0)-f^{*}\bigr)}{\eta T}}_{\text{(I)}}
\;+\;
\underbrace{L\eta\bigl(G^{2}+\sigma^{2}d\bigr)}_{\text{(II)}} .
\tag{1}
\]
Consequently, choosing $\eta =\Theta\!\bigl(1/\sqrt{T}\bigr)$ yields
\[
\frac{1}{T}\sum_{t=0}^{T-1}\E\!\bigl[\norm{\grad f(x_t)}^{2}\bigr]
=O\!\bigl(T^{-1/2}\bigr).
\]
\end{theorem}

\section*{Theoretical Properties}
\begin{itemize}
    \item \textbf{Stability.} The bound (1) shows that as long as the stepsize respects
          $\eta\le 1/L$, the recursion is stable; the additive term (II) is linear in
          the noise variance $\sigma^{2}$.
    \item \textbf{Hyper‑parameter sensitivity.}
          \begin{itemize}
              \item Larger $\eta$ accelerates the deterministic decay (term (I)) but
                    inflates the noise penalty (term (II)).
              \item The optimal $\eta^{\star}= \sqrt{\tfrac{2\bigl(f(x_0)-f^{*}\bigr)}{L\bigl(G^{2}+\sigma^{2}d\bigr)T}}$
                    balances the two terms.
          \end{itemize}
    \item \textbf{Comparison to non‑private SGD.}
          Setting $\sigma=0$ recovers the standard non‑convex SGD bound
          $\frac{2(f(x_0)-f^{*})}{\eta T}+L\eta G^{2}$.
          The privacy‑induced noise adds the extra $L\eta\sigma^{2}d$ term.
\end{itemize}

\section*{Discussion}
The theorem guarantees that even under Gaussian perturbation required for differential
privacy, DP‑SGD makes progress towards stationary points at the same
$O(T^{-1/2})$ rate as vanilla SGD, provided the noise variance is not too large.
In practice, the noise scale $\sigma$ is dictated by the privacy budget
$(\varepsilon,\delta)$; the bound clarifies the explicit trade‑off between
privacy (larger $\sigma$) and optimization accuracy (larger term (II)).

\newpage
\section*{Preliminaries (Definitions and Lemmas)}
\begin{lemma}[Smoothness inequality]\label{lem:smooth}
If $f$ is $L$‑smooth, then for any $x,y$,
\[
f(y)\le f(x)+\langle\grad f(x),y-x\rangle+\frac{L}{2}\norm{y-x}^{2}.
\tag{2}
\]
\end{lemma}
\begin{proof}
Standard consequence of $L$‑Lipschitz gradient; see e.g. Nesterov (2004). \hfill $\square$
\end{proof}

\begin{lemma}[Bound on stochastic gradient variance]\label{lem:var}
Under \textbf{A2} and \textbf{A4},
\[
\E\!\bigl[\norm{g_t-\grad f(x_t)}^{2}\mid x_t\bigr]
\;=\;
\E\!\bigl[\norm{\grad f(x_t;\xi_t)-\grad f(x_t)}^{2}\mid x_t\bigr]
+\E\!\bigl[\norm{\zeta_t}^{2}\bigr]
\;\le\; G^{2}+ \sigma^{2}d .
\tag{3}
\]
\end{lemma}
\begin{proof}
Independence of $\zeta_t$ from $\xi_t$ gives additive variances; $\E\|\zeta_t\|^{2}= \sigma^{2}\E\|Z\|^{2}= \sigma^{2}d$ for $Z\sim\mathcal{N}(0,I_d)$. \hfill $\square$
\end{proof}

\section*{Main Proof (Step‑by‑Step Derivation)}
We prove Theorem~\ref{thm:main}.  All expectations are taken over the randomness
in the mini‑batch $\xi_t$ and the Gaussian noise $\zeta_t$.

\noindent\textbf{Step 1.}  Apply Lemma~\ref{lem:smooth} with $y=x_{t+1}=x_t-\eta g_t$:
\[
f(x_{t+1})\le f(x_t)+\langle\grad f(x_t),-\,\eta g_t\rangle
+\frac{L}{2}\norm{\eta g_t}^{2}.
\tag{4}
\]

\noindent\textbf{Step 2.}  Rearrange (4):
\[
f(x_t)-f(x_{t+1})
\ge \eta\langle\grad f(x_t),g_t\rangle-\frac{L\eta^{2}}{2}\norm{g_t}^{2}.
\tag{5}
\]

\noindent\textbf{Step 3.}  Decompose $g_t$:
\[
g_t = \grad f(x_t) + \underbrace{\bigl(\grad f(x_t;\xi_t)-\grad f(x_t)\bigr)}_{\Delta_t}
       + \zeta_t .
\tag{6}
\]

\noindent\textbf{Step 4.}  Substitute (6) into the inner product of (5):
\[
\langle\grad f(x_t),g_t\rangle
= \norm{\grad f(x_t)}^{2}
  +\langle\grad f(x_t),\Delta_t\rangle
  +\langle\grad f(x_t),\zeta_t\rangle .
\tag{7}
\]

\noindent\textbf{Step 5.}  Substitute (6) into the squared norm term of (5):
\[
\norm{g_t}^{2}
= \norm{\grad f(x_t)}^{2}
  +\norm{\Delta_t}^{2}
  +\norm{\zeta_t}^{2}
  +2\langle\grad f(x_t),\Delta_t\rangle
  +2\langle\grad f(x_t),\zeta_t\rangle
  +2\langle\Delta_t,\zeta_t\rangle .
\tag{8}
\]

\noindent\textbf{Step 6.}  Plug (7) and (8) into (5) and collect terms:
\[
\begin{aligned}
f(x_t)-f(x_{t+1})
&\ge \eta\norm{\grad f(x_t)}^{2}
   +\eta\langle\grad f(x_t),\Delta_t\rangle
   +\eta\langle\grad f(x_t),\zeta_t\rangle \\
&\quad -\frac{L\eta^{2}}{2}\Bigl[
        \norm{\grad f(x_t)}^{2}
        +\norm{\Delta_t}^{2}
        +\norm{\zeta_t}^{2}
        +2\langle\grad f(x_t),\Delta_t\rangle \\
&\qquad\qquad\qquad
        +2\langle\grad f(x_t),\zeta_t\rangle
        +2\langle\Delta_t,\zeta_t\rangle
      \Bigr].
\end{aligned}
\tag{9}
\]

\noindent\textbf{Step 7.}  Take conditional expectation $\E[\cdot\mid x_t]$.  
Using $\E[\Delta_t\mid x_t]=0$ (unbiased stochastic gradient) and
$\E[\zeta_t]=0$, all linear cross‑terms vanish:
\[
\E\!\bigl[\langle\grad f(x_t),\Delta_t\rangle\mid x_t\bigr]=0,
\qquad
\E\!\bigl[\langle\grad f(x_t),\zeta_t\rangle\mid x_t\bigr]=0,
\qquad
\E\!\bigl[\langle\Delta_t,\zeta_t\rangle\mid x_t\bigr]=0 .
\tag{10}
\]

\noindent\textbf{Step 8.}  Apply (10) to (9):
\[
\begin{aligned}
\E\!\bigl[f(x_t)-f(x_{t+1})\mid x_t\bigr]
&\ge \eta\norm{\grad f(x_t)}^{2}
   -\frac{L\eta^{2}}{2}\Bigl[
        \norm{\grad f(x_t)}^{2}
        +\E\!\bigl[\norm{\Delta_t}^{2}\mid x_t\bigr]
        +\E\!\bigl[\norm{\zeta_t}^{2}\bigr]
      \Bigr] .
\end{aligned}
\tag{11}
\]

\noindent\textbf{Step 9.}  Use Lemma~\ref{lem:var} (3) to bound the variance term:
\[
\E\!\bigl[\norm{\Delta_t}^{2}\mid x_t\bigr] + \E\!\bigl[\norm{\zeta_t}^{2}\bigr]
\le G^{2}+\sigma^{2}d .
\tag{12}
\]

\noindent\textbf{Step 10.}  Insert (12) into (11):
\[
\E\!\bigl[f(x_t)-f(x_{t+1})\mid x_t\bigr]
\ge
\Bigl(\eta-\frac{L\eta^{2}}{2}\Bigr)\norm{\grad f(x_t)}^{2}
-\frac{L\eta^{2}}{2}\bigl(G^{2}+\sigma^{2}d\bigr) .
\tag{13}
\]

\noindent\textbf{Step 11.}  Since $\eta\le 1/L$ (Assumption \textbf{A5}),
\[
\eta-\frac{L\eta^{2}}{2}\;\ge\;\frac{\eta}{2}.
\tag{14}
\]

\noindent\textbf{Step 12.}  Combine (13) and (14):
\[
\E\!\bigl[f(x_t)-f(x_{t+1})\mid x_t\bigr]
\ge
\frac{\eta}{2}\norm{\grad f(x_t)}^{2}
-\frac{L\eta^{2}}{2}\bigl(G^{2}+\sigma^{2}d\bigr) .
\tag{15}
\]

\noindent\textbf{Step 13.}  Take full expectation and sum over $t=0,\dots,T-1$:
\[
\sum_{t=0}^{T-1}\E\!\bigl[f(x_t)-f(x_{t+1})\bigr]
\ge
\frac{\eta}{2}\sum_{t=0}^{T-1}\E\!\bigl[\norm{\grad f(x_t)}^{2}\bigr]
-\frac{L\eta^{2}T}{2}\bigl(G^{2}+\sigma^{2}d\bigr) .
\tag{16}
\]

\noindent\textbf{Step 14.}  The left‑hand side telescopes:
\[
\sum_{t=0}^{T-1}\E\!\bigl[f(x_t)-f(x_{t+1})\bigr]
= \E\!\bigl[f(x_0)-f(x_T)\bigr]
\le f(x_0)-f^{*},
\tag{17}
\]
where $f^{*}:=\inf_{x}f(x)$.

\noindent\textbf{Step 15.}  Combine (16) and (17), then divide by $T$:
\[
\frac{1}{T}\sum_{t=0}^{T-1}\E\!\bigl[\norm{\grad f(x_t)}^{2}\bigr]
\le
\frac{2\bigl(f(x_0)-f^{*}\bigr)}{\eta T}
+L\eta\bigl(G^{2}+\sigma^{2}d\bigr) .
\tag{18}
\]
Equation (18) is exactly the bound (1) claimed in Theorem~\ref{thm:main}.
\hfill $\blacksquare$

\section*{Rate Derivation (With Constants)}
From (18), set $\eta = \sqrt{\dfrac{2\bigl(f(x_0)-f^{*}\bigr)}{L\bigl(G^{2}+\sigma^{2}d\bigr)T}}$ (which satisfies $\eta\le 1/L$ for sufficiently large $T$). Plugging this choice yields
\[
\frac{1}{T}\sum_{t=0}^{T-1}\E\!\bigl[\norm{\grad f(x_t)}^{2}\bigr]
\le
2\sqrt{\frac{L\bigl(G^{2}+\sigma^{2}d\bigr)\bigl(f(x_0)-f^{*}\bigr)}\; \frac{1}{\sqrt{T}} .
\tag{19}
\]
Thus the expected squared gradient norm decays as $O(T^{-1/2})$.

\section*{Special Cases}
\begin{itemize}
    \item \textbf{No privacy noise ($\sigma=0$).} Bound (18) reduces to the classical
          non‑convex SGD guarantee with variance $G^{2}$.
    \item \textbf{Purely private update (no stochasticity, $G=0$).} The error term becomes
          $L\eta\sigma^{2}d$, reflecting that privacy alone introduces a variance
          proportional to the ambient dimension $d$.
    \item \textbf{Convex $f$.} If $f$ is additionally convex, the same derivation
          gives a bound on suboptimality $f(\bar x_T)-f^{*}$ rather than gradient norm.
\end{itemize}

\end{document}